\subsection{1.1 Motivation}\label{motivation}

The proliferation of mobile devices has made crowdsourced IoT services
ubiquitous. Consider a pedestrian in a shopping center seeking WiFi
connectivity---their smartphone can connect to hotspots carried by other
shoppers. These moving WiFi hotspots provide service on-the-go,
extending coverage beyond fixed access points. The same principle
applies to vehicle-to-vehicle communication on highways, worker-assisted
sensing in factories, or peer-to-peer data sharing at concerts.

These scenarios share a common challenge: both the service provider and
consumer are moving. The composition that worked moments ago fails now
because the service provider has moved out of range. Unlike static
service composition where services remain available at known locations,
moving IoT services require continuous re-composition as trajectories
diverge.

The core difficulty is spatio-temporal validity. At any given timestep,
a service is valid only if it falls within the consumer's communication
range---the discovery zone. This validity changes moment to moment as
both entities move. Beyond validity, we want services that provide high
quality---measured by channel capacity, which degrades with distance.
The agent must learn to select services that are both valid AND
high-quality, without prior knowledge of which services will be valid.

Traditional service composition approaches fail here because they assume
static services. Even prior DQN-based approaches suffer from three key
limitations: overestimation bias leads to poor action selection when
valid services are scarce, the value-based nature struggles with the
discrete action space of service selection, and most critically, DQN is
reactive---it considers only current positions, ignoring that services
will move.

We need an approach that: (1) learns to select valid services through
interaction, (2) handles discrete action spaces naturally, (3) can
anticipate future states through trajectory encoding, and (4) converges
faster than value-based methods. Advantage Actor-Critic offers all four
capabilities.

\begin{center}\rule{0.5\linewidth}{0.5pt}\end{center}

\hypertarget{background-and-related-work}{}
